\pg{340}

64. O infinito é cognoscivel e conhece-se effectivamente pelo finito.

a) O \emph{infinito é cognoscivel pelo finito}. --- No ente finito, ou limitado, devem distinguitise duas coisas: a \emph{entidade}, ou a perfeição, e a \emph{limitação} a essa entidade, ou a negação de ulterior entidade. A \emph{entidade limitada} não é limitada emquanto é \emph{entidade}, mas emquanto é \emph{limitada}, isto é, emquanto importa a "negação de uma entidade ulterior. Ora a nossa intelligencia, assim como abstrabe de limites determinados e concebe o indefinito, assim tambem pode, sem repugnancia, abstrahir totalmente 'dos limites, e conceber uma entidade sem limites, isto é, uma entidade infinita. Logo o \emph{infinito} é cognoscivel pelo \emph{finito} \footnote{Que \& nossa inelligencia possua a idea do \emph{infinito}, está fora de toda a duvida. Fallamos no infinito, .discutimos acerca do infinito; e ninguem falla numa coisa, ninguem discute ácerca de uma coisa, de que não tém a. idea. A questão refere-se á qualidade d'essa idea Aífirmamos que essa idea é \emph{distincta}, ainda que concedamos que é \emph{imperfeita} e inadequada.\par Os \emph{positivistas} negam-nos a possibilidade de conhecer o \emph{infinito}, dizendo qua a nossa intelligencia só pode conhecer o finito, e que por isso tudo o que estiver além do finito é por nós incognoscivel, --- O erro dos \emph{positivistas} é manifesto. Se o nosso conhecimento devesse --- ou ser sempre e necessariamente \emph{intuitivo}, --- ou comprehender sempre o seu objecto na sua totalidade, --- ou produzil-o, então a intelligencia não poderia conceber o infinito ; porque, sendo ella finita, não pode ter uma intuição directa do infinito, nem comprehendel-o, \emph{e} muito menos produzil-o. Mas é falso que o nosso conhecimento se refira sempre aos objectos, de que se tenha tido uma percepção intuitiva, como é falso que deva sempre comprehender o seu objecto na sua totalidade, ou que deva produzil-o para o conhecer, Logo a intelligencia pode conhecer o infinito, ainda que da um modo finito ou inadequado.}.\footnote{São tres, como dissemos, os elementos, que constituem a idea do \emph{infinito}: a \emph{entidade}, O \emph{limite e a negação} do limite. Ora a nossa intelligencia, considerando os entes finitos, pode conceber. o que é \emph{entidade}. o que é \emph{limite} e o que é negação de todo o limite. Pode conceber o que é \emph{entidade}; porque vê que os entes finitos são alguma coisa e não são nada. Pode conceber o que é \emph{limite}; porque vê que um ente é distincto de outro, um é inferior a outro, e assim percebe que nenhum possue a plenitude da entidade, e que todos são limitados a um genero ou a uma especie, Pode conceber a \emph{negação} de todo o limite; porque, assim como, pela sua força de abstracção, pode afastar de um ente uma determinada limitação, ou imperfeição, percebida noutro ente, assim pode afastar de uma entidade todo o \emph{limite}, toda a imperfeição, e assim conceber, de algum modo, o ente \emph{infinito},\par \emph{Descartes}. \emph{Fabre} e tudos os sequazes do \emph{oníologismo} dizem que não podemos subir do \emph{finito} pare o intinito, porque este não está contido naquelle.\pgn{341}--- Mas tal raciocinio é vicioso. Para que um ente seja conhecido por outro ente, não é necessario que o primeiro esteja contido no segundo \emph{physica} e realmente, mas basta que esteja contido \emph{logicamente}; basta. que o segundo .. tenha com o primeiro uma relação de semelhança ou de dependencia. Ora o finito tem com o infinito uma relação essencial de dependencia e de seme- Ihanea; pois, como veremos, o finito é produzido exclusivamente pelo infinito, e reproduz em si, ainda que de um modo immensamente imperfeito, as perfeições do infinito. Do facto de o infinito não estar \emph{physica} ou \emph{realmente} contido no finito só se deduz que o finito não pode dar-nos uma idea adequada do infinito, o que concedemos e sustentamos ; mas não se deduz que o finito não possa dar-nos nenhuma idea, nem sequer inadequada, do infinito; porque. para isso, basta afastar da entidade ou perfeição finita todo e qualquer limite. Será um conhecimento imperfeito do ente perfeitissimo; mas será sempre conhecimento.}

